\chapter{Локализация смены ранга нейтринной опоры}
\label{ch:v6-spectral-transition-rank-change-localization-gate}

\section{Различающий вопрос}

Квадратичный ковариант
\begin{equation}
 W_\nu(H)=\widetilde H\widetilde H^\dagger
\end{equation}
регулярен во всём пространстве конфигураций и меняет ранг с нуля на один
ровно при переходе от $H=0$ к $H\neq0$. Однако существование такого
оператора ещё не означает, что область смены ранга локализуется на уже
построенном вихревом дефекте $Q,T,B$.

Нулевая гипотеза настоящего гейта утверждает, что радиальное сцепление
$M_{300}$ само подавляет норму Хиггса в ядре вихря до нуля. Стоп-критерий
--- строго положительный точечный минимум $|H|^2$ для всякой допустимой
неотрицательной радиальной амплитуды.

Это новая проверка, а не повторение портального аудита. Предыдущий гейт
установил нулевой коэффициент оператора
$\Tr(Q^2)H^\dagger H$, но сохранил радиальный член $T|H|^2$. Здесь
проверяется, достаточно ли именно этого сохранившегося члена для
локального восстановления слабой симметрии.

\section{Точечное сокращение полного квадрата}

Единый квадрат кривизны Тома V имеет вид
\begin{equation}
 V_H(X,H)=\frac13\left\{
 \Tr(I_3-X^\dagger X)^2+
 \Tr(XX^\dagger-|H|^2I_3)^2+
 3(|H|^2-1)^2\right\}.
 \label{eq:v6-rank-localization-parent-potential}
\end{equation}
Положим
\begin{equation}
 XX^\dagger=TR,
 \qquad T\geq0,
 \qquad \Tr R=1,
 \qquad s=|H|^2\geq0.
\end{equation}
После удаления членов, не зависящих от $s$, получаем
\begin{equation}
 V_{\rm loc}(s;T)=\frac13\left[6s^2-(2T+6)s\right]+\mathrm{const}.
 \label{eq:v6-rank-localization-reduced-potential}
\end{equation}
Его единственный минимум на $s\geq0$ равен
\begin{equation}
 s_*(T)=\frac{T}{6}+\frac12.
 \label{eq:v6-rank-localization-current-minimum}
\end{equation}
Следовательно,
\begin{equation}
 s_*(T)\geq\frac12>0
 \qquad\text{при всех }T\geq0.
\end{equation}

Знак радиального сцепления оказался решающим: положительная амплитуда
$T$ не подавляет, а повышает предпочтительную норму Хиггса. Поэтому даже
идеализированная точка $T=0$ не является точкой восстановления $H=0$.

\section{Оптимистический контроль единичного ядрового члена}

В старом гессианном аудите $M_{300}$ уже рассматривался единичный
дефектный импульс $V_k=|H|^2$. Чтобы не пропустить этот наиболее дешёвый
канал, введём контрольный параметр $0\leq k\leq1$:
\begin{equation}
 V_{\rm loc}^{(k)}(s;T)
 =\frac13\left[6s^2-(2T+6-3k)s\right]+mathrm{const}.
\end{equation}
Тогда
\begin{equation}
 s_*(T;k)=\max\left\{0,
 \frac{T}{6}+\frac12-\frac{k}{4}\right\}.
 \label{eq:v6-rank-localization-optimistic-minimum}
\end{equation}
Даже на наиболее благоприятной границе $T=0$, $k=1$ остаётся
\begin{equation}
 s_*\geq\frac14>0.
\end{equation}
Для достижения $H=0$ потребовалась бы сила
\begin{equation}
 k\geq k_{\rm crit}(T)=2+\frac23T,
 \label{eq:v6-rank-localization-critical-strength}
\end{equation}
которой нет в текущем родителе. Контроль $k=1$ не объявляется новым
членом модели: он лишь показывает, что даже уже встречавшаяся
единичная нормировка недостаточна.

\section{Следствие для спектрального перехода}

Для слабого дублета
\begin{equation}
 \Tr W_\nu(H)=|H|^2=s.
\end{equation}
Поэтому при найденном минимуме
\begin{equation}
 \rank W_\nu=1
\end{equation}
как вне вихря, так и в его ядре. Текущий дефект может деформировать
амплитуду $W_\nu$, но не создаёт границу ранговых страт $0\leftrightarrow1$.

Это согласуется с общей теорией конденсатов на струнах: второе поле
обнуляется или конденсируется в ядре только при наличии смешанного
взаимодействия с требуемым знаком и достаточной величиной; одна топология
струнообразующего поля такого условия не обеспечивает
\cite{rank-localization-witten1985,rank-localization-achucarro2000,
rank-localization-forgacs2017}.

\begin{theorem}[Радиальный мост не локализует смену ранга]
В существующем квадрате кривизны $M_{300}$ точечный минимум нормы слабого
дублета равен $|H|^2=T/6+1/2$ и потому строго положителен при $T\geq0$.
Даже оптимистический единичный ядровый член оставляет
$|H|^2\geq1/4$. Следовательно, вихрь $Q,T,B$ не локализует переход
$\rank W_\nu:0\to1$.
\end{theorem}

\begin{nogocriterion}[Запрет ручного ядрового подавления]
Нельзя назначить $k\geq k_{\rm crit}(T)$ или ненулевой коэффициент
$\Tr(Q^2)H^\dagger H$ ради получения $H=0$ в ядре. Такой шаг вводит
новое взаимодействие после знания желаемого результата. Допустим только
его вывод из представленного коннектора, общего квадрата кривизны или
другого заранее определённого родительского принципа.
\end{nogocriterion}

\section{Следующий различающий тест}

Отрицательный ответ закрывает радиальную интерпретацию существующего
моста, но не все внутренние углы $M_{300}$. Следующий гейт о радиальном
коннекторе вихря должен классифицировать представленные эквивариантные стрелки между
семейным дефектным сектором и слабым дублетом. Успех требует ненулевого
бифундаментального пространства, чей единый квадрат автоматически даёт
ядровый член правильного знака. Нулевое пространство интертвинеров
закроет скрытый коннектор и отделит текущий родитель от новой модели.
Его файловый идентификатор:
\begin{sloppypar}
\path{version6_spectral_transition_radial_bridge_vortex_connector_gate}.
\end{sloppypar}

\begin{protocol}
\begin{enumerate}
 \item регулярный носитель смены ранга $W_\nu$: \textbf{сохранён};
 \item радиальное сцепление $T|H|^2$: \textbf{учтено};
 \item точечный минимум текущего родителя:
 \textbf{$|H|^2=T/6+1/2$};
 \item нижняя граница при $T\geq0$: \textbf{$1/2$};
 \item нижняя граница при оптимистическом $0\leq k\leq1$:
 \textbf{$1/4$};
 \item критическая сила подавления:
 \textbf{$k_{\rm crit}=2+2T/3$};
 \item локализованный переход ранга $0\to1$:
 \textbf{не получен};
 \item скрытый портал формы: \textbf{не найден};
 \item физическое замыкание: \textbf{не достигнуто}.
\end{enumerate}
\end{protocol}

Машинный сертификат сохранён в
\path{s2t_v6_spectral_transition_rank_change_localization_gate.py} и
\path{s2t_v6_spectral_transition_rank_change_localization_gate_results.json}.

\begin{thebibliography}{9}
\bibitem{rank-localization-witten1985}
E.~Witten,
\emph{Superconducting Strings},
Nuclear Physics B 249 (1985), 557--592.

\bibitem{rank-localization-achucarro2000}
A.~Ach\'ucarro, T.~Vachaspati,
\emph{Semilocal and Electroweak Strings},
Physics Reports 327 (2000), 347--426; arXiv:hep-ph/9904229.

\bibitem{rank-localization-forgacs2017}
P.~Forg\'acs, A.~Luk\'acs,
\emph{Stabilization of Semilocal Strings by Dark Scalar Condensates},
Physical Review D 95 (2017), 035003; arXiv:1612.03151.
\end{thebibliography}